\originalchapter{光力与激光冷却}
\originalsection{散射力与偶极势}
吸收行波光子带来平均动量$\hbar\mathbf k$, 自发发射在对称角分布下平均为零, 但有非零方差. 因而$\mathbf F_{\rm sc}=\hbar\mathbf k\Gamma\rho_{ee}$. 最强单向辐射压力为$\hbar k\Gamma/2$. 用第\ref{ch:open}章的稳态布居, 两束弱相互干扰的反向红失谐光给
\[F(v)=\frac{\hbar k\Gamma s}{2}\left[\frac1{1+s+[2(\Delta-kv)/\Gamma]^2}-\frac1{1+s+[2(\Delta+kv)/\Gamma]^2}\right].\]
对$kv\ll\Gamma$展开得$F=-\beta v$, 红失谐$\Delta<0$时
\[\beta=-\frac{8\hbar k^2s\Delta/\Gamma}{[1+s+(2\Delta/\Gamma)^2]^2}>0.\]
速度朝哪束光运动, 就更接近哪束的共振, 从而受到阻力.

远失谐光的保守作用为$U=-\alpha(\omega)I/(2\epsilon_0c)$. 靠近一个跃迁,$\Delta=\omega_L-\omega_0$, 基态势为$U\simeq\hbar\Omega^2/(4\Delta)$; 红失谐$U<0$吸向高强度, 蓝失谐排斥高强度. 散射率$\Gamma\Omega^2/(4\Delta^2)$, 故$R_{\rm sc}/|U|=\Gamma/(\hbar|\Delta|)$. 增大失谐可减小固定势深下的散射, 但需要更高功率且应使用多能级极化率.

\originalsection{动量扩散与冷却极限}
随机吸收与发射使动量扩散. 一维Fokker--Planck方程$\partial_tf=(\beta/m)\partial_p(pf)+D_p\partial_p^2f$的Gaussian稳态满足$\langle p^2\rangle=mD_p/\beta$, 所以$k_BT=D_p/\beta$. 理想闭合两能级、弱饱和模型计入自发反冲后得到
\[T=\frac{\hbar\Gamma}{4k_B}\frac{1+(2\Delta/\Gamma)^2}{2|\Delta|/\Gamma},\qquad T_{\min}=\frac{\hbar\Gamma}{2k_B}\quad(\Delta=-\Gamma/2).\]
摩擦不能单独决定温度; 它必须与扩散比较. 真实多能级偏振梯度可实现亚Doppler冷却, 利用不同光移势和光抽运产生Sisyphus机制. 单次反冲能$E_R=(\hbar k)^2/(2m)$为更低尺度, 但不是所有方案的绝对温度下限.

\originalsection{Zeeman减速器、MOT与磁阱}
原子束被持续散射减速后, Doppler位移不断改变. Zeeman减速器令$\Delta-kv(z)-\delta_Z(z)\simeq0$, 用空间磁场补偿速度变化. 对恒定减速度$a$, $v(z)=\sqrt{v_0^2-2az}$, 磁场轮廓由共振条件给出. 工作加速度应留出余量低于最大辐射压力加速度, 否则原子会脱离共振.

磁光阱用四极磁场使左右位置的Zeeman频移相反, 再用相应圆偏振使离中心的原子更吸收指向中心的光. 在小位移小速度下$F=-\kappa_xx-\beta v$. 磁场提供位置辨识而不是直接以磁力捕获所有原子. 高密度时再吸收荧光带来斥力, 形成密度限制.

纯磁阱势为$U=m_Fg_F\mu_B|B|$. 低场寻求态可被磁场最小值捕获, 高场寻求态在静态源外磁场中通常不能形成三维稳定极大值. 四极阱零点可引起非绝热自旋翻转; 偏置场阱或光阱可以避免相应损失.
